% HEADLINE: Firmware \& Embedded Linux Platform Engineer -- Zynq UltraScale+ / ARM SoCs, kernel, drivers, real-time Ethernet
% Persona B (Staff firmware / embedded Linux platform IC), 2 pages max. Tightened from base Content/WorkExperience.tex into pyramid form:
%   Summary first; bullets re-sorted by argument; every number traced to a Bullet Bank ID; unreconciled numbers replaced by \textbf{⚠️ ID}.
% Arguments: (1) owns the platform, not a feature; (2) works at the kernel / driver / hardware boundary; (3) delivers under release pressure and integrates new stacks.
% Removed from base: "10 times (520ms to 50ms)" (W-01, ⚠️ R1 — no W-01/W-03/W-04 used here), "pipelining and dynamic caching" (W-04-adjacent, ⚠️ R7),
%   "9+" projects (⚠️ R11), "reducing manual QA time by 80%" (I-04, ⚠️ R10), "96.7%" (I-01 → 96.75 %, ANPR moved to Publications).
% Ortiq omitted (default B = Big Tech / silicon-vendor IC loop). For OEM AT / Helsing / contractor builds add one "Other" line: Founder & CTO, Ortiq (2026).
% Lab titles: Kazan as on the cv. RaccoonLab = "Principal Engineer" [stated 2026-10-29] on every build (overrides the old Big-Tech downgrade rule; Bullet Bank §4, ⚠️ R25 verification).

\section{Summary}
\cvitem{}{Platform firmware engineer who owns the Linux stack -- bootloader to application -- for three industrial sensor product families on Zynq UltraScale+ and ARM64 SoCs. Brings kernel and driver work (V4L2, I2C, DMA), real-time Ethernet integration, and the release discipline that shipped FW 1.3.0.}   % W-06 (3 families) · W-09 · W-07 · W-15 · W-12 · W-05 (qualitative, no %) — 44 words

\section{Experience}

\cventry{Nov 2023 -- Present}{Embedded Linux Software Developer}{Wenglor MEL GmbH}{Unterschleißheim, Germany}{}{   % F-01 ⚠️ R4: cv title kept until R4 is ticked
\begin{itemize}
  \item Owned the Yocto distribution and Jenkins CI/CD (build, HIL test, deploy) for the G4 (Zynq UltraScale+) and B60 camera lines; migrated them to kernel 6.1 / 5.15.   % W-09 + W-10 (Arg 1) — scope first (Benchmark P3: Apple/Big Tech read product + silicon scope); lines + SoC from W-10; kernel 6.1 / 5.15
  \item Led the software-licensing infrastructure end to end (WIBU CodeMeter: architecture, firmware integration, rollout); Wenglor now earns revenue by selling sensor feature licenses.   % W-13 (Arg 1) — [stated 2026-10-29]; no revenue/SKU/headcount numbers (⚠️ R33 R34 R35)
  \item Architected a modular nginx/GenICam-style Ethernet vision server now reused across 3 product families without frame-rate loss.   % W-06 (Arg 1) — 3 families
  \item Introduced read-only rootfs and A/B firmware updates for Cyber Resilience Act readiness; proposed container isolation for user apps.   % W-11 (Arg 1) — 2 shipped mechanisms + 1 proposal
  \item Brought up Sony IMX296/IMX567 MIPI CSI-2 sensors on AMLogic A311D and Zynq UltraScale+ (V4L2 sub-devices, capture FSMs); root-caused a B60 I2C bus-busy fault in the kernel and TCP drops on the Zynq sensor (Wireshark).   % W-07 (2 sensors, 2 SoCs) + W-15 + W-16 (Arg 2) — merged so the role keeps ≤ 7 bullets (Style Guide §1 depth rule; Benchmark P4)
  \item Carried \textbf{⚠️ W-05} of the ShapeDrive G4 firmware 1.3.0 load and fixed its release-blocking memory-corruption crash.   % W-05 (Arg 3) — ⚠️ R6 (~80 %) include-decision pending
  \item Prototyped EtherNet/IP on Hilscher CiFX/netX90, evaluated PROFINET and EtherCAT.   % W-12 (3 protocols) — Arg 3; W-13 moved to its own top-3 bullet
\end{itemize}}

\cventry{Jan 2023 -- Sep 2023}{OEM Process Manager}{Yandex, Smart Devices -- Factory Team}{Moscow, Russia}{}{
\begin{itemize}
  \item Built one Python toolchain (OpenHTF, ADB, U-Boot, Ansible/Salt) automating 4 factory workflows: flashing, SN/MAC, HDCP keys, QA.   % Y-02 — 4 workflows
  \item Re-shaped QA/QC/AQL on the Shenzhen OEM lines and introduced metrics-driven defect tracking.   % Y-04 — qualitative; scope = OEM lines
\end{itemize}}

\cventry{Feb 2021 -- Jan 2023}{Embedded Development Lead}{Kazan Scientific Center, Quantum Optics Lab}{Kazan, Russia}{}{
\begin{itemize}
  \item Designed the high-speed Spatial Light Modulator control protocol and public SDK: bare-metal Cortex-M firmware plus Linux drivers over RNDIS/USB and MIPI CSI-2.   % K-01 — scope: several optical experiments
  \item Set up reproducible builds and containerised CI (GitHub Actions, Docker); published the open-source Python SDK (holoeye\_embedded) with Sphinx docs.   % K-02 + K-03 — public repo github.com/Ehsan2754/HOLOEYE-PLUTO-2-EMBEDDED
\end{itemize}}

\cventry{Aug 2019 -- Jan 2021}{Principal Engineer}{RaccoonLab (Innopolis)}{Innopolis, Russia}{}{   % [stated 2026-10-29] org + title; dates ⚠️ R24
\begin{itemize}
  \item Developed autonomous-drone CAN modules end to end: Altium hardware, C/C++ FreeRTOS firmware, UAVCAN/DroneCAN/Cyphal protocol layer, bench tooling.   % I-06 (Arg 2, lead) — no count: ⚠️ R19 open, non-numeric form
  \item Took the modules into production: ran the ramp-up, estimated cost per batch and qualified drop-in substitutes for components that went scarce mid-run.   % I-07 (Arg 3) — no units/duration: ⚠️ R20 R21 open, non-numeric form
  \item Integrated the nodes with PX4/ArduPilot autopilots and automated firmware builds and HIL simulation in GitHub Actions.   % I-03 + I-10 (+ HIL half of I-05) — 2 autopilot stacks; I-04 (−80 % QA) still omitted: ⚠️ R10
  \item Co-invented a registered visual-positioning system for autonomous multirotor landing (patent RU 2022610258).   % I-02 / P-05 — 1 patent
\end{itemize}}

\cventry{2013 -- 2019}{Embedded \& PCB Developer}{Freelance and FUM Robotics Lab}{Mashhad, Iran}{}{
\begin{itemize}
  \item Shipped AVR firmware and PCBs for a dental X-ray exposure controller, a wireless kiln thermometer and an A/C servicing machine.   % E-02 — 3 named boards
  \item Built STM32/FreeRTOS localisation nodes over LoRaWAN, UWB (DW1000) and GNSS for the City of Mashhad (FUM Robotics Lab).   % E-02 + E-09 — 3 radio stacks; customer from long cv L124
\end{itemize}}
